%! Exponential and Logarithmic Equations and Inequalities — Unit 5, Lesson 5.6 — Exit Ticket (Answer Key)
\documentclass[10pt]{article}
\usepackage{precalculus-article}
\usepackage{precalculus-key}   % pulls in -boxes; do NOT also load -boxes
\newcommand{\TallMath}[1]{$\displaystyle #1\rule[-1.4em]{0pt}{3.2em}$}

\begin{document}

\pageheader{Unit 5, Lesson 5.6}{Exit Ticket --- Answer Key}
\namedateperiod

\vspace{0.14in}

\begin{enumerate}[leftmargin=1.8em, itemsep=18pt]

\item Solve $5^x = 80$. Show all steps. Give your answer in exact form and as a decimal
      approximation to 3 decimal places.

\vspace{6pt}
Step 1: \ansline{Take log of both sides: $\log(5^x) = \log 80$}

\vspace{6pt}
Step 2: \ansline{Power rule: $x \cdot \log 5 = \log 80$; divide: $x = \log 80 / \log 5$}

\vspace{6pt}
Exact: $x = $ \ans{$\log 80 / \log 5$} \quad (equivalently $\log_5 80$)

\vspace{6pt}
Decimal: $x \approx $ \ans{2.723}

\vspace{10pt}

\item Solve $\log_4(x + 2) = 3$.

\vspace{6pt}
Work:

\vspace{6pt}
\ansline{Convert to exponential form: $x + 2 = 4^3 = 64$; subtract 2: $x = 62$.}

\vspace{8pt}
Solution: $x = $ \ans{62}

\vspace{8pt}
Check for extraneous solutions: Does this value make the argument of the logarithm positive?

\vspace{6pt}
Argument value: $x + 2 = $ \ans{64} \qquad Is this $> 0$? \ans{Yes}

\vspace{6pt}
Is the solution valid? \ans{Yes} \quad Justify: \ansline{$64 > 0$, so the argument is in the domain of $\log_4$. No extraneous solution.}

\end{enumerate}

\end{document}
